\documentclass[10pt]{article}

    

\def\Type{LessonDJ}
\def\TypeNum{Lesson 0}
\def\TypeTitle{Modeling Change and Uncertainty\\ Notes to Instructors }
\def\Semester{Fall 2021} %Not Used 
\def\Course{MATH 1190 \& 1210}
\def\Targets{\bf Intro the course + Pre-Test!}

\usepackage{preambleDJ} 

\begin{document}
\pagestyle{otherpages}
%\thispagestyle{empty}
\thispagestyle{firstpage}


\vs{.85}

If you see the phrase \mod in the Notes To Instructors, that means you should feel free to adapt to your personality and style.  Most of the mini-lectures embedded in the handouts this semester will be something you can modify to suit your personality, voice and perspective.   \\\\
When you see the phrase {\em Core problems} this will indicate the core problems of the handout. So if you are not sure what to focus on, or are concerned about running out of time, focus on the Core Problems and skip the others.\\\\
This first day of class is always unique as we are trying to do some important things to get things off to a good start. Thus goals for today include several around setting the right tone/vibe for learning the rest of the course. I'm going to try to avoid spending time on administrative details that are in the syllabus (i.e. how much each exam is worth, etc.) and focus on other things.\\\\
\lt  {\bf Intro the course + Pre-Test!} Today is a day unlike any other in the course! We want to start the process of:
\itemI{
\item Building a learning community that values a growth mindset, a sense of belonging,  a purpose mindset that connects learning to other important areas of their life,  and struggle as a normal part of learning.
\item Introducing ourselves and the Learning Assistant as real human beings!
\item Introducing the Key Ideas of the course and active learning methods we will be using.
\item Guiding students towards what is necessary for a successful week 1.
\item Take the Pre-Test in last $50$ minutes of class.
}
~\\
Typically, a description of each LT that will be assessed will be provided above.  Note that  {\bf Intro the course + Pre-Test!}  is not a learning target (LT) that will be tested.  Sometimes, LTs span more than one lesson. That means in each lesson will have {\em Lesson Goals} for that specific lesson that feed into the larger Learning Targets.
\\\\
\lg 
\enum{
\item Begin to create an atmosphere where students will talk to each other, to you, and to the Learning Assistant about mathematics and things that happen outside of the classroom.
\enum{
\item Humanize yourself by telling students who you are and some things about yourself and why you teach. 
\item The Learning Assistant (LA) introduces himself/herself and talks about "What it's like taking a math class at UVA."
\item Have students talk to each other about something. This could involve mathematics or something not related to mathematics.
}
\item Set course expectations.
\enum{
\item Growth mindset.  Everyone can and will grow in how they learn mathematics. We will be working to help students adopt a growth mindset over the course of the semester. It's important to start building this culture today.
\item Interactive in-class experience.  We will be using active learning techniques including students frequently working together.
\item Normalize Struggle.  We will expect a lot! Students will feel challenged in ways that they might not have experienced before. For many students, this will be the first time they have ever struggled in math.  This is normal. With hard work and effective practices, this feeling will go away.  {\em If you can share a time where you struggled/failed with something, that can be very helpful for students}.  
}
\item Talk about course structure
\enum{
\item Pre-class work: Combination of reading texts and/or watching video(s) and answering a few questions. Last question will always be {\em ``Is there anything in the pre-class assignment that wasn't clear or that you didn't understand? Was there anything you found especially interesting?"}
\item In-class work will be tailored by feedback from pre-class work.
\item Post-class work will be WebAssign HW.
\item Application Days will help us see how math ideas we're wrestling with show up in the world around us.
\item Growth Based group vs Standard Group NEED MORE HERE
\item Tools we will use in the course: Poll Everywhere to incentivize in-class learning and sometimes Desmos. If you have other tools you think you will use Discord, Slack, Geogebra, Mathematica etc), today is a good day to use them for something.   
}
\item Begin to communicate to students ``I care about you and your learning. We're going to grow together." I'm also going to collect student profiles  that I will use throughout the semester to highlight students each time we are together.  Feel free to duplicate and \href{https://drive.google.com/drive/u/2/folders/1YQsjblmECOi7H3XIXbhu6sbbKhNxoCDx}{\underline{use this or not}}.
\item Do some mathematics! After today, the students will be able to:
\itemI{
\item Explain what the course goal is: { \em To Use Mathematical Tools to Describe Change and Uncertainty}.  
\item Create a model of the amount of CO$_2$ in the atmosphere from a data set and use it to predict a future amount of CO$_2$.
\item Reflect on process of creating model and suggest other possible models/strategies.
}
\item Take Pre-Test. Please set aside about $50$ minutes for this.
}



\feed
\itemC{
\item Since there is no pre-class activity, we do not have any feedback from prior semesters.
}

 \hand
 \itemI{
 \item  {\em Introduction} The first page of class is time for the instructor and LA to introduce themselves and the course.
\item The Keeling Curve exercise is designed to have students start to think mathematically about  real-world data and the ideas of change and uncertainty.  $75$ minute classes have roughly $25$ minutes to get through pages $1$ and problems $1$ and $2$
\item Pre-test. The last 50 minutes of class are devoted to a pre-test.
 }

\core Page 1 (i.e. {\bf\blue{Welcome!}}), \# $1, 2$ and {\bf \blue{Key Questions for the semester}}.\\
Extra, if time:  problems \#$3-4$ will set up students for success on our first Application Day in a couple of weeks.




\newpage
{\bf Before Class} \mod Think about the tone you want to set as students come to class. I'm going to be playing some songs because I will be starting Student Of The Day next class and a part of that will be playing a student's favorite songs.  So I'm starting with that today, except they are my songs instead of the students'.  {\em You should not feel any obligation to do what I'm doing here.}
\itemC{
\item Carmin Burina
\item Take it to the limit (Eagles)
\item Welcome to the Jungle (Guns N Roses)
\item I will survive (Gloria Gaynor)
}
~\\
\bblue{Welcome!} \mod \\
 \mpageT{.5}{.5}[\hs{.1}]{
 {\em A bit about me...}
  \nti{
  \itemC{
  \item  Humanize yourself: I'm going to talk about where I'm from, some of my academic history and work experience and then mention a few personal things (hobbies, family, the fact that I root for Duke in basketball, etc.).  
  \item Normalize struggle.  I'm also going to tell students that everyone hits a wall at some point.  I will describe what that was like for me in grad school when I hit an academic wall for the first time in my life. Some will have this experience in Math 1320 and so I'm going to try to normalize struggle. I will point out that struggle is normal and part of the learning process. If you are not struggling you probably aren't learning! It's normal and effort and helpful learning practices (working with others, talking to you as instructor, going to P2L, MCLC, etc.) will help them grow as learners and lead to these feelings going away.
  \item Growth Mindset. As part of my conversation the relationship between struggle and learning, I will talk about how we can grow in how we learn and that I expect everyone to grow in this way this fall.  We will want to repeat this message several times over the course of the semester.
  }
 }
 }{
{ \em Now, about you...}
 \nti{
 I'm going to have them introduce themselves to each other in groups  using the prompt "Who gave you your name? Why that name? Do you know the ethnic origin of your name? Do you have nicknames?\\\\
 This will be their first experience of talking with each other and starting to build an environment that I hope will pay off when we do mathematics.
 }
 }\vs{.1}
   {\em We've got a Learning Assistant!}  I will have my LA introduce themselves.  They will also talk about "What it's like taking a math class at UVA" from the student perspective.\\\\
 {\em Here's where we're headed...}\\
\mpageT{.5}{.5}[\hs{.1}]{
I.C.E
 \nti{
 Here's where I'm going to talk about some details of the course. In particular, I will talk about my teaching philosophy: "You learn math best by you doing math, not by watching me do math." I will also talk about the fact that the course is flipped, the ICE paradigm and what this entails:
 \itemC{
 \item  I= Inform via videos before class
 \item C= Confirm via prep quizzes before class
 \item E= Extend knowledge during and after class.
 }  It might be a good idea  to show them the  Canvas site and where they can find lesson-by-lesson requirements to do before class. I'm going to have students fill out profile sheets electronically before coming to class. But this is also a great time to have students do this (see NTI on first page for link).
 }
}{
Growth Based Assessments (if time)
\nti{
If we are running short on time, I will skip this as there will be plenty of time over next few lessons to talk about this.  If I have time, I will explain:
\itemC{
\item We're working hard to improve their learning experience.  
\item We're implementing something called ``Growth Based Assessments" where we hope to foster an environment where students grow in their learning.
\item We have identified $29$ Learning Targets for the semester. Student will be expected to become proficient in order to get credit for that Learning Target.  In particular,  Final grade will be determined by how many proficiencies are earned.
\item "Proficient" means PUT DESCRIPTOR HERE. So we give students multiple chances to become proficient.  
\item Students will be placed in one of two categories: Growth Based or Standard for each exam.
\item Standard students will be graded the usual way.
\item We promise that at the end of the semester if one group has a lower GPA than the other, the lower group will be curved to meet the higher group.
\item Details are provided in the syllabus.
}
}
}

\newpage

{\bf What's the story so far?}. \mod \\
\mpageT{.4}{.6}[\hs{.1}]{
The Keeling Curve is the record of atmospheric CO2 in parts per million (ppm). This data set is named after Charles Keeling who started collecting this data at Mauna Loa Observatory, starting in 1958.  A sample of atmospheric CO$_2$ is shown below in five-year intervals from starting in 1965 through 2010.  Note that ``year" means ``years since 1965."  For example, $t=10$ represents the year 1975.
}{
\graphicsC{keelingDataMMA}{3.5}
}

\begin{center}
\begin{tabular}{|c|c|c|c|c|c|c|c|c|c|c|}
\hline
year & 0 & 5 & 10 & 15 & 20 & 25 & 30 & 35 & 40 & 45 \\
\hline
CO$_2$ level & 320 & 325 & 331 & 338 & 345 & 354 & 360 & 369 & 379 & 389 \\
\hline
\end{tabular}
\end{center}

{\bf How much  CO$_2$ is in the atmosphere in the year 2025?}
\enum{
\item \pe Let's start with some intuition.  What's your best guess for the amount of CO$_2$ in the atmosphere in the year 2025?  
\nti{
The intention of this question is to introduce tools we will be using this semester (Poll Everywhere) and to have students begin to think about using mathematics in a context. I expect many students students to eyeball the graph and say something in the ballpark of $410$ ppm. It's fine if they assume it's a linear function and extend the graph in some fashion. It's also fine to assume it's an increasing function and extend the graph in some manner. I'm going to use Poll Everywhere to deliver this question because it is a tool we will use throughout the semester. But feel free to put people in groups. 
}
\vs{1}
\item \pe What decisions did you make and/or questions did you have when you made your best guess? 
\nti{
The intention of this question is to begin to draw out of students things that we can connect to mathematics and to encourage students to ask questions. Again, I'm using Poll Everywhere but feel free to put people in groups. Possible responses:
\itemI{
\item I needed to figure out which value of $t$ corresponds to $2025$ ($t=60$)
\item I looked at the graph and it looks like it continues to around $410$ when $t=60$
}
Possible prompts to help students think about this question:
\itemI{
\item What assumptions did you make on first question?
\item What value of $t$ corresponds to the year $2025$?
\item What might we assume about how the amount of CO$_2$ will  grow?
}
}
\newpage
\item  (Extra, if time) Devise at least two strategies to refine your estimate and/or increase your confidence in your estimate.   You might find \url{https://bit.ly/desmosKeeling1} helpful.
\nti{
The intention of this question is to get their minds thinking about many reasonable approaches to modeling the data and making predictions. We will ask them to use some of the ideas that they generate today in our first Application Day concerning the Keeling Curve.\\\\
Possible answers:
\itemC{
\item Collecting more data points would be helpful.
\item Instead of eyeballing the graph and guessing, maybe graph some kind of a function?
\item I assumed data was linear, but it might be bending up.
}
Some possible prompts:
\itemC{
\item Can you think of a kind of function that you have seen before who's graph might look something like this?
\item Try graphing a function in Desmos that seems to mimic the data set in Desmos.
\item Would collecting more data help things? Hurt things?
}
}
\item (Extra, if time) What other information might be helpful to further  improve your estimate and/or build confidence in your estimate? What mathematical tools might be helpful? What remaining questions do you have (you probably have many questions!)?
\nti{
I'm assuming there will be many answers here. Our job is to encourage people think about what mathematics might be useful and how that relates to Change or Uncertainty. Possible strategies
\itemI{
\item Draw a line on the graph
\item Use Desmos to create a linear function and graph it
\item Use Desmos to create a non-linear function and graph it.
\item Use the table of data to measure how much things are increasing every $5$ years
}
Possible things to build confidence:
\itemI{
\item More data than every 5 years and hopefully more recent data.
\item Is there evidence  CO$_2$ will continue to grow, level out, or begin to decrease?
\item Can we estimate probability that if CO$_2$ increases over one year that it will continue to increase the next year? What's the probability that if it increases, it will stay the same the next year? Decrease the next year?
}
Actual amount of CO$_2$ in January of 2024:  $422.94$ ppm.  January 2025 will likely be higher.  We will develop mathematical tools and revisit this scenario in a couple of weeks.
}
}

\bblue{Key Questions for the semester}
\nti{
\enum{
\item How can we use mathematics to describe and understand change?
\item How can we use mathematics to navigate uncertainty?
}
For example, 
\itemC{
\item A biologist might be interested in population variation over time.
\item A medical researcher might be interested in a model of blood pressure as a function of body
weight, or concentration of a drug in the bloodstream over time.
\item A business executive might study the demand for a product as a function of its price or, perhaps,
as it relates to the amount spent on marketing.
\item An environmental scientist might be interested in modeling the level of a toxin in a lake over
time.
\item A data scientist might be interested in fitting a regression line to a data set to make predictions
}
}

\end{document}